\SourceLeaf{89}
\Article{1}
To explain magnetic phenomena, we postulate two magnetic fluids. With the physicists we call one northern and the other southern. We suppose that the elements of either fluid attract those of the other, whereas two elements of the same fluid repel one another, and that both actions vary inversely as the square of the distance. It will be seen below that our own observations fully confirm this law.

These fluids do not appear by themselves, but only in association with the weighable particles of bodies capable of magnetism. Their actions manifest themselves by tending to set these particles in motion, or by preventing or altering the motion that other forces acting on them, such as gravity, would produce.

Thus the action of a given quantity of magnetic fluid on a given quantity of either the same or the other fluid at a given distance will be comparable to a given motive force, that is, to the action of a given accelerating force on a given mass. Since we can know the magnetic fluids themselves only through the effects they produce, those very effects must serve to measure them.

To give this measure a definite meaning, we must first establish units for three kinds of quantity, namely distance, weighable mass and accelerating force. For the third we may take gravity at the place of observation. If this choice is less satisfactory, a unit of time must be introduced as well. We then take the accelerating force to be $=1$ if, in one unit of time, it changes the velocity of a body moving in its direction by an amount equal to the unit of velocity.

\SourceLeaf{90}
With these conventions, the unit quantity of northern fluid is the quantity whose repulsive action on an equal quantity at a distance $=1$ is equivalent to a motive force $=1$, that is, to the action of an accelerating force $=1$ on a mass $=1$. The same applies to the unit quantity of southern fluid. In this determination, both the acting fluid and the fluid acted upon must clearly be conceived as concentrated at physical points. We must further suppose that the attraction between given quantities of unlike fluids at a given distance equals the repulsion between respectively equal quantities of like fluids. The action of a quantity $m$ of northern magnetic fluid on a quantity $m'$ of the same fluid at a distance $r$, supposing both to be concentrated at physical points, is therefore expressed by $\frac{mm'}{rr}$. Equivalently it is a motive force $=\frac{mm'}{rr}$ acting in the direction from the former quantity to the latter. Clearly this formula holds generally if we take the quantity of southern fluid as negative, as we shall always understand from now on. A negative value of the force $\frac{mm'}{rr}$ then indicates attraction in place of repulsion.

Thus if equal quantities of northern and southern fluid are present together at a physical point, no action at all will arise from them. If the quantities are unequal, only the excess of one fluid comes into consideration. We shall call it \textit{free} magnetism, positive or negative.

\Article{2}
To these fundamental assumptions we must add another, confirmed by experience on all sides. Every body containing magnetic fluids always contains an equal quantity of each. Indeed, experience shows that the assumption must extend even to the individual parts of such a body, however small, provided our senses can distinguish them. But as we noted at the end of the preceding article, an action can exist only insofar as the fluids are in some measure separated. We must therefore suppose that they separate across intervals so small that we cannot measure them.

We must therefore conceive of a body capable of magnetism as an aggregate of innumerable particles, each containing a certain quantity of northern magnetic fluid and an equal quantity of southern fluid. The fluids are either uniformly mixed, in which case the magnetism is latent, or have undergone some greater or lesser separation, in which case the magnetism is developed. This separation can never lead to a transfer of fluid from \SourceLeaf{91}one particle to another. It makes no difference whether we suppose a greater separation to arise from a greater quantity of the fluids becoming free or from a greater interval between them. Given the magnitude of the separation, its direction must clearly be considered as well. Depending on whether the directions in the different particles of the body agree or oppose one another, the total action at points outside the body may be greater or smaller.

However free magnetism may be distributed inside the body, a more general theorem always allows us to substitute, according to a definite law, another distribution on the body's surface alone, exactly equivalent to the former with respect to the forces acting outward. An element of magnetic fluid placed anywhere outside the body will thus experience exactly the same attraction or repulsion from the actual distribution of magnetism inside the body as from the fictitious one on its surface. We may extend the same fiction to \textit{two} bodies that act on each other through the free magnetism developed in them. For each body the fictitious surface distribution can then take the place of the actual internal distribution. Only in this way can we give a true meaning to the common manner of speaking that attributes only northern magnetism to one end of a magnetic needle, for example, and southern magnetism to the other. This expression clearly does not agree with the fundamental principle stated above, which other phenomena imperatively demand. Let it suffice to note this here in passing. The theorem itself is unnecessary for our present purpose, and we shall discuss it more fully on another occasion.

\Article{3}
The \textit{magnetic state} of a body consists in the way free magnetism is distributed among its individual particles. With respect to the mutability of this state we observe an essential difference among bodies capable of magnetism. In some, for example soft iron, the slightest force changes the state at once, and the former state returns when the force ceases. In others, especially hardened steel, the force must reach a certain intensity before it can produce a perceptible change in the magnetic state. Once the force ceases, the body either stays in the state it has acquired or at least does not fully return to its former state. Thus in bodies of the first kind the molecules of the magnetic fluid always arrange themselves in perfect equilibrium under the forces they exert on one another and those arising from external causes, or at least they hardly depart perceptibly from such equilibrium. In bodies of the second kind, by contrast, the \SourceLeaf{92}magnetic state can persist even without perfect equilibrium among those forces, provided stronger external forces are kept away. Though the cause of this phenomenon is unknown, we may picture it as if the weighable parts of a body of the second kind offered an obstacle like friction to the motion of the magnetic fluids associated with them. In soft iron this resistance is either absent or at least very slight.

These two cases require entirely different treatment in theoretical investigations, but in the present treatise we shall speak only of bodies of the second kind. In the experiments to be discussed, the stability of the magnetic state in each of the bodies used will be a fundamental assumption. We must therefore take great care during the experiments to keep other bodies that might alter this state from coming too close.

There is, however, one cause of change to which bodies of the second kind are also subject, namely heat. Experience shows that a body's magnetic state varies with its temperature and that increased heat weakens the intensity of its magnetism. Provided the body has not been heated excessively, its former magnetic state returns when its former temperature is restored. This dependence must be determined by suitable experiments. If operations belonging to the same experiment have been carried out at unequal temperatures, the first step is to reduce them to the same temperature.
